% ============================================================================
% Глава 3: Архитектура и реализация
% ============================================================================

\section{АРХИТЕКТУРА И РЕАЛИЗАЦИЯ}

\subsection{Структура на данните}

Системата за управление на автопарк с IoT интеграция използва три основни колекции в MongoDB, които естествено моделират бизнес логиката на логистичната платформа. Структурата е проектирана с акцент върху ефективността на четене и гъвкавостта при еволюция на системата.

\subsubsection{Колекция vehicles}

Колекцията съхранява информация за превозните средства в автопарка. Всеки документ представлява отделно превозно средство и съдържа технически характеристики, статус и история на поддръжката.

\begin{lstlisting}[language=JavaScript,caption=Примерен документ от колекцията vehicles]
{
  "_id": "V001",
  "make": "Mercedes-Benz",
  "model": "Actros",
  "year": 2020,
  "load_capacity": 25,
  "fuel_type": "diesel",
  "status": "active",
  "maintenance_history": [
    {
      "date": "2023-01-15T00:00:00.000Z",
      "description": "Oil change"
    },
    {
      "date": "2023-06-20T00:00:00.000Z",
      "description": "Tire replacement"
    }
  ]
}
\end{lstlisting}

\subsubsection{Колекция shipments}

Тази колекция съдържа информация за товарните пратки, които се транспортират от превозните средства. Документите включват маршрутна информация, теглителни характеристики, статус на доставката, както и времеви рамки (\texttt{start\_time} и \texttt{end\_time}), които дефинират пряката връзка с телеметричните данни на съответното превозно средство в този времеви период.

\begin{lstlisting}[language=JavaScript,caption=Примерен документ от колекцията shipments]
{
  "_id": "S001",
  "origin": "Sofia",
  "destination": "Plovdiv",
  "weight": 15,
  "priority": "high",
  "assigned_vehicle_id": "V001",
  "status": "delivered",
  "start_time": "2023-10-01T06:00:00.000Z",
  "end_time": "2023-10-01T12:00:00.000Z",
  "estimated_arrival": "2023-10-01T10:00:00.000Z"
}
\end{lstlisting}

\subsubsection{Колекция telemetry}

Колекцията съхранява сензорни данни в реално време от IoT устройствата, монтирани на превозните средства. Това е най-динамичната колекция с висока честота на записване.

\begin{lstlisting}[language=JavaScript,caption=Примерен документ от колекцията telemetry]
{
  "vehicle_id": "V001",
  "timestamp": "2023-10-01T08:00:00.000Z",
  "gps": {
    "lat": 42.6977,
    "lng": 23.3219
  },
  "metrics": {
    "speed": 80,
    "fuel_level": 85.5,
    "engine_temp": 90
  }
}
\end{lstlisting}

Тази структура демонстрира гъвкавостта на документо-ориентирания модел - полето `readings` в сензорите позволява съхранение на времеви серии директно в документа, а референциите към превозни средства осигуряват ефективни връзки \cite{json_schema}.

\subsection{Изпълнявани операции}

\subsubsection{CRUD операции}

Системата реализира пълния набор от CRUD операции върху трите колекции:

\textbf{Създаване (Create):}
\begin{itemize}
    \item Добавяне на ново превозно средство с основни характеристики и празен списък сензори
    \item Регистрация на сензор с начални параметри и празен масив за показания
    \item Създаване на запис за поддръжка с планирана дата и статус "pending"
\end{itemize}

\textbf{Четене (Read):}
\begin{itemize}
    \item Филтриране по статус и тип гориво: \texttt{db.vehicles.find(\{status: "active", fuel\_type: "diesel"\})}
    \item Търсене по регулярен израз в регистрационни номера: \texttt{db.vehicles.find(\{license\_plate: \{\$regex: "CA", \$options: "i"\}\})}
    \item Филтриране по диапазон на капацитет: \texttt{db.vehicles.find(\{capacity: \{\$gte: 20000, \$lte: 30000\}\})}
    \item Комбинирани условия с OR: \texttt{db.maintenance.find(\{\$or: [\{status: "completed"\}, \{status: "in\_progress"\}]\})}
\end{itemize}

\textbf{Актуализиране (Update):}
\begin{itemize}
    \item Добавяне на ново показание към сензор: \texttt{db.sensors.updateOne(\{\_id: ObjectId("...")\}, \{\$push: \{readings: \{value: 90.0, timestamp: new Date(), status: "normal"\}\}\})}
    \item Актуализиране на локация на превозно средство: \texttt{db.vehicles.updateOne(\{\_id: ObjectId("...")\}, \{\$set: \{location: \{latitude: 42.7, longitude: 23.3, last\_update: new Date()\}\}\})}
    \item Увеличаване на общите разходи за поддръжка: \texttt{db.maintenance.updateOne(\{\_id: ObjectId("...")\}, \{\$inc: \{total\_cost: 50.00\}\})}
\end{itemize}

\textbf{Изтриване (Delete):}
\begin{itemize}
    \item Премахване на неактивни сензори: \texttt{db.sensors.deleteMany(\{is\_active: false\})}
    \item Изтриване на конкретен запис за поддръжка: \texttt{db.maintenance.deleteOne(\{\_id: ObjectId("...")\}\})}
\end{itemize}

\subsubsection{Агрегиращи операции}

За анализ на данни е реализирана агрегираща заявка за намиране на превозни средства с най-високи разходи за поддръжка \cite{mongodb_definitive}:

\begin{lstlisting}[language=JavaScript]
db.maintenance.aggregate([
  { $match: { status: 'completed' } },
  {
    $group: {
      _id: '$vehicle_id',
      totalMaintenanceCost: { $sum: '$total_cost' },
      maintenanceCount: { $sum: 1 },
      averageCost: { $avg: '$total_cost' }
    }
  },
  { $sort: { totalMaintenanceCost: -1 } },
  { $limit: 5 },
  {
    $lookup: {
      from: 'vehicles',
      localField: '_id',
      foreignField: '_id',
      as: 'vehicle'
    }
  },
  { $unwind: '$vehicle' },
  {
    $project: {
      licensePlate: '$vehicle.license_plate',
      model: '$vehicle.model',
      totalMaintenanceCost: 1,
      maintenanceCount: 1,
      averageCost: { $round: ['$averageCost', 2] }
    }
  }
])
\end{lstlisting}

Тази заявка демонстрира използването на операторите \texttt{\$match}, \texttt{\$group}, \texttt{\$sort}, \texttt{\$limit}, \texttt{\$lookup} и \texttt{\$project} за комплексен анализ на данни от две колекции.

\subsection{Сигурност и контрол на достъпа}

В системата е реализиран Role-Based Access Control (RBAC) чрез функционалностите на MongoDB. Създадени са следните роли за управление на достъпа:
\begin{itemize}
    \item \textbf{DataAnalyst}: Роля с права само за четене (\texttt{find}) върху всички колекции (\texttt{vehicles}, \texttt{shipments}, \texttt{telemetry}). Използва се за аналитични нужди без риск от модификация на данните.
    \item \textbf{FleetManager}: Роля с пълни CRUD права (\texttt{find}, \texttt{insert}, \texttt{update}, \texttt{remove}) върху колекциите \texttt{vehicles} и \texttt{shipments}, но само с право на четене върху \texttt{telemetry}.
    \item \textbf{LogisticsApp}: Системна роля, предназначена за IoT устройствата, с права единствено за запис (\texttt{insert}) в колекцията \texttt{telemetry}.
\end{itemize}

Ролите са назначени на съответните потребители (\texttt{data\_analyst}, \texttt{fleet\_manager}, \texttt{logistics\_app}), гарантирайки принципа на най-малката привилегия (Principle of Least Privilege).

\subsection{Тестване и валидация}

Системата е покрита с автоматизирани unit тестове (чрез Jest), които валидират целостта на данните и логиката на връзките. Специално внимание е обърнато на негативните тестове за проверка на валидността на времевите рамки на пратките и ограниченията за достъп. 

Примерни логове от симулирани невалидни и неоторизирани заявки, прихванати по време на тестването:

\begin{lstlisting}[language=bash,caption=Логове от отхвърлени заявки и неоторизиран достъп]
[Error Logged] Invalid request: Shipment S999 missing start or end time
[Error Logged] Invalid request: start_time cannot be after end_time for Shipment S999
[AccessControl Error] User with role 'DataAnalyst' denied insert on vehicles.
  -> MongoServerError: not authorized on fleetTracking to execute command { insert: "vehicles" }
[AccessControl Error] User with role 'FleetManager' denied insert on telemetry.
  -> MongoServerError: not authorized on fleetTracking to execute command { insert: "telemetry" }
[AccessControl Error] User with role 'LogisticsApp' denied find on vehicles.
  -> MongoServerError: not authorized on fleetTracking to execute command { find: "vehicles" }
\end{lstlisting}

Архитектурата на системата осигурява ефективно управление на IoT данни с висока степен на гъвкавост и производителност, като същевременно поддържа целостта и консистентността на информацията \cite{martin2017clean}.